% Source: https://artofproblemsolving.com/community/c3h2670574_some_actually_good_proving_problems
% Stable identity: aops-topic-2670574-q17
Take a circle with centre $O$ and radius $r$, and mark six points on it at $60^\circ$ intervals; they form a regular hexagon. Each side subtends $60^\circ$ at $O$, so each triangle formed by $O$ and a side is isosceles with apex angle $60^\circ$, hence equilateral, and every side of the hexagon has length $r$.

\begin{asy}
usepackage("amsmath");
size(3cm);
defaultpen(fontsize(10pt));
draw(unitcircle,linewidth(0.8));pair O=(0,0);dot(O);
pair[] P;for(int i=0;i<6;++i)P[i]=dir(60*i);
draw(P[0]--P[1]--P[2]--P[3]--P[4]--P[5]--cycle);
draw(O--P[0],dashed);draw(O--P[1],dashed);
label("$O$",O,SW);label("$r$",(O+P[0])/2,S);label("$r$",(O+P[1])/2,W);label("$r$",(P[0]+P[1])/2,SW,fontsize(8pt));
draw(arc(O,0.2,0,60));label("$60^\circ$",0.2*dir(30),NE,fontsize(7pt));
\end{asy}

Each side of the hexagon is a straight segment joining the ends of one of the six arcs into which the circle is divided, so it is shorter than that arc. Adding over the six arcs,
\[6r<2\pi r,\qquad\text{so}\qquad\pi>3.\ \blacksquare\]
